\section{坐标测量中的事件选择}
\subsection{从光锥得到变换系数}
惯性系的时空均匀性要求：同一个坐标差采用同一变换，与起点位置、起始时刻无关。坐标关系因而具有仿射形式；将原点对齐后，变换为线性。先把原点对齐，沿$x$的变换可写$x'=a x+bct$、$ct'=d x+ect$。$S'$原点在$S$中满足$x=\beta ct$，因此$b=-a\beta$。

分别代入向右、向左的两条原点光线。$x=ct$要求$x'=ct'$，所以$a+b=d+e$；$x=-ct$要求$x'=-ct'$，所以$-a+b=d-e$。两式相加、相减给$d=b=-a\beta$、$e=a$。至此光锥不变只确定了整体比例$a$，尚未确定其数值。用相对性原理要求逆变换在交换两系后只有$\beta$变号，正逆矩阵相乘便给$a^2(1-\beta^2)=1$。连续到$\beta=0$的恒等变换，选$a=\gamma>0$。

这一推导也说明为什么单独要求光速不变还不能任意省略比例系数。若同时允许改变坐标单位，光锥形状不变但标尺标度也会改变；本书要求两系采用相同物理单位和互易的标准钟尺。

横向坐标不变可由对称性理解：沿$x$的运动没有选出横向某一方向。设横向统一缩放为$k(\beta)$，正逆变换要求$k(\beta)k(-\beta)=1$，空间反射对称给两者相等，连续正根为1。这排除了把横向尺寸也按$\gamma$缩短的做法。

\subsection{同一事件、同一物体与同一时刻}
坐标题有三种常见约束，各自消去不同未知量。两条世界线相遇要求同一事件的$t,x$都相同；跟踪同一物体是在其世界线上取两事件；同时测长度则是在某一系固定$t$后取不同物体的事件。只有把题设约束先转为这些关系，代入变换才有明确意义。

固定一个钟的两个事件满足$x'=\text{常数}$，因而$\Delta x=V\Delta t$。代入时间变换得$\Delta t'=\gamma\Delta t(1-\beta^2)=\Delta t/\gamma$。若改成在$S$同时读取两只钟，便有$\Delta t=0$而不是$\Delta x=V\Delta t$；相同的式子转而给$\Delta t'=-\gamma V\Delta x/c^2$。两种情境都由同一变换式产生，不能把它们当作两个可随意替换的“钟慢规则”。

\Diagram{fig:m1pairings}{沿同一钟的两事件与沿同一等时线的两事件，是不同配对}{\draw[axis](-.2,0)--(4.4,0)node[right]{$x$};\draw[axis](0,0)--(0,4.5)node[above]{$ct$};\draw[curve](0,0)--(2.4,4);\draw[curve](1,0)--(3.4,4);\fill(.6,1)circle(1.5pt)node[left]{$A$};\fill(1.8,3)circle(1.5pt)node[left]{$B$};\fill(2.8,3)circle(1.5pt)node[right]{$C$};\draw[guide](1.8,3)--(2.8,3);}
图\ref{fig:m1pairings}中$A,B$在同一钟世界线上，$B,C$在同一实验室等时线上。$AB$适合求钟经历的时间，$BC$适合求实验室钟距。若问题还给定钟面被设过零，最后须在坐标时间上加钟的设置常数，不能让这个常数反过来改变世界线。

\subsection{四钟模型的完整事件坐标}
例\ref{ex:m1rods}可把所有坐标集中写成无量纲表。定义$g=\gamma(1+\beta)$、$a=(1-1/\gamma)/\beta$，坐标单位分别为$t_0,l$。$A$钟读$t/t_0$，$B$钟读$t/t_0-a$；$A'$钟读$t'/t_0$，$B'$钟读$t'/t_0+a$。表中每行都代表一个事件，列之间只做坐标变换。
\begin{center}\small
\begin{tabular}{ccccc}\toprule
事件&$t/t_0$&$x/l$&$t'/t_0$&$x'/l$\\\midrule
$M$&0&0&0&0\\
$N$&$a$&1&$-a$&1\\
$P$&1&1&$1/g$&$1/g$\\
$R$&$g$&$g$&1&1\\
$Q$&$2g-1$&1&$\gamma(2g-1-\beta)$&$\gamma[1-\beta(2g-1)]$\\
$T$&$2\gamma$&$2\gamma\beta$&2&0\\
$U$&$2g$&0&$2\gamma g$&$-2\gamma\beta g$\\\bottomrule
\end{tabular}
\end{center}
$P$行使用$\gamma(1-\beta)=1/g$；$R$行则由源系最简单的$(1,1)$变来。$N$行的$t'_N=-t_N$是本题两杆等固有长和特定归零事件共同造成的结果，换成不等长两杆后一般不再成立。

表中的$Q,T$两行也重现例\ref{ex:m1rods}已求出的返回次序及临界速度$\beta=1/\sqrt5$。几何上比较$x_T=2\gamma\beta l$与固定$B$位置$l$即可判断：反射光向左走，先遇到更右侧的对象。$\beta=.6$时$A'$在右，故先$T$后$Q$。

这些次序针对有因果联系的返回光线，同一系或另一系都会给出相同先后。局部归零后的两只不同钟却可能读出反向大小，因此“表上较大的读数”不能单独代替事件次序。$\beta\to0$时$a\simeq\beta/2\to0$，公式连续趋于静止往返，但正好$\beta=0$时两杆端点一直重合，所谓第二次相遇归零事件需要另作规定。

\section{四维量的几何和角分布的测度}
\subsection{任意方向的纯推动}
沿单位矢量$\hat{\vct V}$分解空间矢量$\vct r=\vct r_\parallel+\vct r_\perp$。对平行分量使用一维变换，垂直分量保持不变，得到
\[
 ct'=\gamma\left(ct-\frac{\vct V\cdot\vct r}{c}\right),\quad
 \vct r'=\vct r+\left[\frac{\gamma-1}{V^2}(\vct V\cdot\vct r)-\gamma t\right]\vct V.
\]
$V=0$的极限取恒等映射。对任意四矢量，只需把$ct$换成其时间分量、把$\vct r$换成空间分量。四动量的能量变换因而是$E'=\gamma(E-\vct V\cdot\vct p)$。这里$\vct V$是坐标系速度，$\vct p$是被研究粒子的动量，二者没有理由总是平行。

洛伦兹标量在任意允许的惯性变换下都保持数值，如固有时、静质量和两个四矢量的缩并。某一分量在一次特殊变换中恰好不变，还不足以判定它是标量；光子能量就可能在特定方向和推动速度下不变，换一种变换后又改变。

类时四矢量的空间分量可在某个系全为零。非零类光四矢量的闵可夫斯基平方为零，时间与空间分量都不能通过有限惯性变换消去；真空光子因而没有静止系。类空四矢量可以把时间分量变为零，却不能把空间分量全消去。这些分类取决于闵可夫斯基内积；时空图的纸面长度不表示四矢量模。

\subsection{单位角、单位立体角与环带}
轴对称分布中，单位立体角强度$I(\theta)$与极角环带内功率不同。环带固体角$\dd\Omega=2\pi\sin\theta\dd\theta$，故$\dd P/\dd\theta=2\pi I(\theta)\sin\theta$。即使$I(\theta)$在$\theta=0$最大，有限环带功率的密度也因$\sin\theta$在轴上为零而不在那里最大。画角分布前须说明纵轴是哪种密度。

\Diagram{fig:m1solidring}{极角环带宽度相同，立体角仍随$\sin\theta$改变}{\draw[curve](0,0)circle(1.7);\draw[axis](0,0)--(2.7,0)node[right]{轴};\draw[curve](0,0)--(1.472,.85);\draw[curve](0,0)--(1.302,1.093);\draw[thick](1.472,.85)arc(30:40:1.7);\node[above]at(.4,2){$\dd\Omega=2\pi\sin\theta\dd\theta$};}
图\ref{fig:m1solidring}以截面短弧表示球面环带。极坐标曲线的半径表示所选的角密度，积分时仍须乘立体角雅可比。运动源的功率份额可先选定一个系中的角域，再变换边界计算，往往能省去完整角分布的积分。

更一般地，源系辐射不必各向同性。设单位固有时发出的总能量为$P_0$，总$x$向动量率为$\dot p'_x$。整个发射小过程的四动量变换给$\dd E=\gamma(\dd E'+\beta c\dd p'_x)$，而发射时间$\dd t=\gamma\dd\tau$，于是实验室总发射功率为
\[
 P=P_0+\beta c\dot p'_x.
\]
各向同性源有$\dot p'_x=0$，所以总功率不变；向后单束光有$\dot p'_x=-P_0/c$，所以$P=P_0(1-\beta)$。两种结果都由同一四动量式得到。总功率不变需要源系净辐射动量为零，各向同性发射满足这一条件。

\subsection{密度变换中的超平面与通量}
粒子数守恒并不表示不同参考系中的数密度相等。计数所用的体积必须处在该系的同时面上；同一批粒子在两个同时面上被截取的位置不同。对共同四速度为$U^\mu$的稀疏粒子流，定义数流$N^\mu=n_0U^\mu$。在实验室中，
\[
 N^\mu=(cn,\vct j_N)=n_0\gamma_v(c,\vct v),\qquad
 n=\gamma_vn_0,\quad \vct j_N=n\vct v.
\]
穿过时空小面元的粒子数为$\dd\mathcal N=N^\mu\dd\Sigma_\mu/c$。实验室等时小体积有$\dd\Sigma_\mu=(\dd V,0,0,0)$，给出$n\dd V$；跨越固定空间面的面元则给出粒子通量。密度与流量是同一四维对象的不同分量，变换时不能仅替换长度而保持时间率不变。

无热运动的粒子束还具有能量动量流$T^{\mu\nu}=n_0mU^\mu U^\nu$。其分量为
\[
 T^{00}=n_0m\gamma_v^2c^2=u,\quad
 cT^{0x}=uv_x,\quad T^{xx}=n_0m\gamma_v^2v_x^2.
\]
最后一项是单位面积、单位时间跨面的$x$动量；即使束在自身系没有压力，运动系也测得定向动量通量。对存在随机热运动的气体，须对每个粒子的速度分布求和，不能再把整个系统视为只有一个共同$U^\mu$的冷束。
\Diagram{fig:m1currentcut}{同一粒子流由不同同时面截取；斜线不代表粒子轨迹的偏折}{\draw[axis](0,0)--(4.5,0)node[right]{$x$};\draw[axis](0,0)--(0,3.7)node[above]{$ct$};\foreach\a in{.4,1.2,2,2.8}{\draw[curve](\a,0)--({\a+.7},3.2);}\draw[guide](.1,1.1)--(3.8,1.1)node[right]{$t=\text{常量}$};\draw[densely dotted](.1,1.6)--(3.8,2.7)node[right]{$t'=\text{常量}$};}
图\ref{fig:m1currentcut}的两条截线各自连接一组不同事件。密度变换应同时指明计数的同时面和跨面通量，两者由同一四流联系。静质量密度$n m$、总能量密度$u$和$u/c^2$分别包含不同的洛伦兹因子，不能互作同义量。
